% ============================================================
%  ch00：摘要 + §1 引言（原文 pp. 1--6）
%  底本：arXiv:1602.06282v3
% ============================================================

\section*{摘要}
\addcontentsline{toc}{section}{摘要}

我们构造了一个函子，从 $p$ 进域上光滑刚性解析簇 $X$ 上的 $p$ 进平展局部系统范畴，
映到其「到 $\BdR$ 的基变换」上带可积联络的向量丛范畴；这可视作迈向人们寻求的
$p$ 进 Riemann--Hilbert 对应的第一步。作为推论，我们得到连通刚性解析簇上
$p$ 进局部系统的如下刚性定理：若该局部系统在一点处的茎（视为 $p$ 进 Galois 表示）
在 Fontaine 意义下是 de Rham 的，则它在每一点处的茎都是 de Rham 的。
在此过程中，我们还建立了 $p$ 进 Simpson 对应的一些基本性质。最后，
我们给出这些结果在志村簇上的一个应用。

\section{引言}\label{sec:intro}

我们先陈述本理论的一些应用。

\begin{thm}\label{thm:1.1}
设 $X$ 是数域 $E$ 上的几何连通代数簇，$L$ 是 $X$ 上的 $p$ 进平展局部系统。
假设存在闭点 $x\in|X|$，使得 $L$ 在 $x$ 的某个几何点 $\bar x$ 处的茎 $\bar L_x$
（视为 $x$ 的剩余域的 $p$ 进 Galois 表示）在 Fontaine--Mazur 意义下是几何的
（即几乎处处非分歧，且在 $p$ 处为 de Rham），
则对每个闭点 $y\in|X|$，$L$ 在 $y$ 处的茎 $\bar L_y$
（视为 $y$ 的剩余域的 Galois 表示）也是几何的。
\end{thm}

在继续之前，我们先作几点说明。

\begin{rmk}\label{rmk:1.1}
回忆：由代数簇的平展上同调得到的 $p$ 进表示都是几何的，而 Fontaine--Mazur 猜想预测其逆。
鉴于上述定理，下述 Fontaine--Mazur 猜想的相对版本看来是合理的。
\end{rmk}

\begin{conj*}
设 $L$ 是数域 $E$ 上几何连通代数簇 $X$ 上的平展 $\Qp$ 局部系统，$\eta$ 是 $X$ 的泛点。
若对某个 $x\in|X|$，$\bar L_x$ 是几何的，则存在 $\eta$ 上的代数簇 $Y$，
使得 $L_{\bar\eta}$ 在 Tate 扭的意义下作为 $Y_{\bar\eta}$ 的平展上同调的子商出现。
\end{conj*}

注意：若 $L$ 的单值群是交换的，则该猜想成立。事实上，此时几何单值群必有限，
于是在取有限平展覆盖之后 $L$ 成为几何常值的，即某个（$E$ 的有限扩域上的）
交换 Galois 表示的拉回。由于交换情形下的 Fontaine--Mazur 猜想是已知的（例如见 \cite{He}），
上述猜想在此情形下也成立。

\begin{rmk}\label{rmk:1.2}
不妨将上述定理与 Deligne 的下述结果 \cite{DM}（该文称之为 Principle B）作比较：
设 $S$ 是光滑连通复代数簇，$f:X\to S$ 是紧合光滑态射。设 $\{t_s\}_{s\in S}$
是由 $S$ 参数化的一族 Hodge 类（即相对上同调 $R^{2d}f_*\Q(d)$ 的一个整体截面，
它在每个 $s\in S$ 处限制为 $H^{2d}(X_s,\Q(d))$ 中的 Hodge 类）。
若它在某一点处是绝对 Hodge 的，则 $t_s$ 对所有 $s\in S$ 都是绝对 Hodge 的。
\end{rmk}

下面是定理 \ref{thm:1.1} 的一个具体应用，它源于与蓝凯文（K.-W. Lan）的讨论。
设 $(G,X)$ 是志村数据。对（充分小的）开紧子群 $K\subset G(\Af)$，令
\[
  \Sh_K(G,X)=G(\Q)\backslash\bigl(X\times G(\Af)\bigr)/K
\]
为相应的志村簇。设 $V$ 是 $G$ 的 $\Q$ 有理表示，且它在 $Z_G^s$ 上平凡。
这里 $Z_G^s$ 是 $G$ 的中心中最大的、在 $\R$ 上分裂的各向异性子环面。\zy{%
原文作 ``the largest anisotropic subtorus in the center of $G$ that is $\R$-split''，
字面上「各向异性」与「$\R$-分裂」相抵牾。按志村簇的通常约定，此处应指中心的最大
$\R$-分裂子环面；原文表述存疑，暂按字面译出，待核。}
众所周知，$V$ 在 $\Sh_K(G,X)$ 上诱导一个 Betti 局部系统 $L_V$。此外，
典范模型理论给出 $\Sh_K(G,X)$ 的一个定义在反射域 $E\subset\C$ 上的模型（仍记作同一记号），
并且对素数 $p$ 的选取，给出 $\Sh_K(G,X)$ 上的一个 $p$ 进平展局部系统 $L_{V,p}$
（参见 \cite[§III.6]{Mi1}）。把定理 \ref{thm:1.1} 应用于志村簇上的特殊点，便得到下述定理。

\begin{thm}\label{thm:1.2}
保持上述记号与假设。则对 $\Sh_K(G,X)$ 的每个闭点 $x$，
茎 $(L_{V,p})_{\bar x}$（视为 $\Gal(\bar E/E(x))$ 的 Galois 表示）
在 Fontaine--Mazur 意义下是几何的。
\end{thm}

证明及相关结果见第 4.2 节。

\begin{rmk}\label{rmk:1.3}
当 $(G,X)$ 为阿贝尔型（再加上少量额外假设）时，定理 \ref{thm:1.2} 是已知的：
此时 $\Sh_K(G,X)$ 参数化某些阿贝尔动机，而 $L_{V,p}$ 正是它们的 $p$ 进实现。
在一般情形，该定理为 Deligne 的预期——志村簇（在少量额外假设下）应当是某些动机的模空间——
提供了证据（尤其结合注 \ref{rmk:1.1} 中的猜想来看）。
\end{rmk}

现在我们转向本文的主要理论。设 $k$ 是 $\Qp$ 的有限扩域，$\bar k$ 是 $k$ 的代数闭包，
$\khb$ 是其完备化。定理 \ref{thm:1.1} 的主要构件是下述定理。

\begin{thm}\label{thm:1.3}
设 $X$ 是 $k$ 上的几何连通刚性解析簇，$L$ 是平展 site $X_{\text{\'et}}$ 上的 $\Qp$ 局部系统。
若存在 $X$ 的经典点 $x$，使得 $L$ 在 $x$ 的某个几何点 $\bar x$ 处的茎 $\bar L_x$
（视为 $x$ 的剩余域的 $p$ 进表示）是 de Rham 的，
则对 $X$ 的每个经典点 $y$，茎 $\bar L_y$ 都是 de Rham 的。
此外，$\bar L_y$ 与 $\bar L_x$ 具有相同的 Hodge--Tate 权（计重数）。
\end{thm}

\begin{rmk}\label{rmk:1.4}
不难看出，定理 \ref{thm:1.3} 的陈述无论是晶体版本还是半稳定版本都不成立。
事实上，可以考察由 $\Qp$ 上光滑连通刚性解析簇参数化的一族阿贝尔簇，
使得其中大多数纤维在 $p$ 处有好归约，但某些纤维没有半稳定归约；
于是平凡局部系统 $\Qp$ 在该族上的高阶顺像就在底簇上给出所需的反例。
另一方面，人们预期定理 \ref{thm:1.3} 的潜半稳定版本成立：
即若该族在一个经典点处是半稳定的，则在 $X$ 的某个有限平展扩张之后，
甚至在基域的某个有限扩张之后，它在每个经典点处都变成半稳定的。
\end{rmk}

\begin{rmk}\label{rmk:1.5}
$X$ 上的 $\Qp$ 局部系统可以看作由 $X$ 参数化的一族几何 $p$ 进表示。
不妨将定理 \ref{thm:1.3} 与 $p$ 进表示的算术族的相关结果 \cite{BC,Be} 作比较。
后者说的是：若一个算术族的 $p$ 进表示在底簇的一个 Zariski 稠密子集上为 de Rham，
且相应的 Hodge--Tate 权集合一致有界，则该族是 de Rham 的。
此外，特征簇上出现的 $p$ 进表示族表明，「一致有界」这一条件是不可缺少的。
因此，定理 \ref{thm:1.3} 展示了算术族与几何族之间的二分：后者具有惊人的刚性。
\end{rmk}

\begin{rmk}\label{rmk:1.6}
对素数 $\ell\neq p$，Kisin 证明了 $\Zl$ 局部系统的一个刚性定理：
在剩余特征为 $p$ 的非阿基米德域上的有限型概形上，$\Zl$ 局部系统在 $p$ 进拓扑下局部常值 \cite{Ki}。
\end{rmk}

利用刚性解析簇的奇点消解（参见 \cite{BM}），只需对光滑簇证明定理 \ref{thm:1.3}。
在此情形，我们将从一个版本的 $p$ 进 Riemann--Hilbert 对应推出它；下面就来解释这一对应。

设 $\BdR$ 表示 Fontaine 的 de Rham 周期环。在 $\Xk$ 上考虑 $\Qp$-代数层
$\OX\widehat{\otimes}\BdR$（精确定义见 §3.1）；它从 $\BdR$ 上的滤过
继承一个滤过，并从 $\OX$ 上的导子继承一个 $\BdR$-线性导子。
环化空间 $(\Xk,\OX\widehat{\otimes}\BdR)$ 记作 $\mathcal X$，
可以把它看作「$\Xk$ 到 $\BdRp$ 的某个典范提升的泛纤维」。
我们有下述定理，它可视为迈向 $p$ 进簇上人们长期寻求的 Riemann--Hilbert 对应的第一步。

\begin{thm}\label{thm:1.4}
（完整而精确的陈述见定理 3.8。）设 $X$ 是 $k$ 上的光滑刚性解析簇。
则存在一个张量函子 $\mathrm{RH}$，从 $X_{\text{\'et}}$ 上的 $\Qp$ 局部系统范畴，
映到 $\mathcal X$ 上的向量丛范畴（即某些有限局部自由的 $\OX\widehat{\otimes}\BdR$-模），
这些向量丛带有 $\Gal(\bar k/k)$ 的半线性作用，以及一个滤过和一个满足 Griffiths 横截性的可积联络。
函子 $\mathrm{RH}$ 与任意态射的拉回相容，并且在某些条件下与光滑紧合态射的顺像相容。
\end{thm}

\begin{rmk}\label{rmk:1.7}
在 \cite{Be} 中，算术 Galois 表示族的 de Rham 周期层被定义为 $X$ 的结构层与 $\BdR$
在基域上的完备张量积。然而，由于「典范嵌入 $\khb\to\BdR$」不存在，
$\Xk$ 上的层 $\OX\widehat{\otimes}\BdR$ 必须用稍微迂回的方式定义。
另一个关键区别是：在几何情形，向量丛上还带有额外结构，即可积联络。
正如下文将看到的，这使得从 $\OX\widehat{\otimes}\BdR$-模到 $\OX$-模的过渡
比算术族的情形更简单、更漂亮。
\end{rmk}

我们把定理 \ref{thm:1.4} 视为 $p$ 进平展局部系统的几何 Riemann--Hilbert 对应。
现在设 $\phi:\Xk\to X$ 是自然投影。由于对 $X$ 上的局部系统 $L$，
向量丛 $\mathrm{RH}(L)$ 带有 $\Gal(\bar k/k)$ 的作用，故（非正式地）可以定义
\[
  D^i_{\dR}(L)=H^i\bigl(\Gal(\bar k/k),\phi_*\mathrm{RH}(L)\bigr),\qquad i\ge 0.
\]
更精确的定义见 §3.2。于是定理 \ref{thm:1.3}（对光滑 $X$）
可由下述定理推出，后者可视为算术 Riemann--Hilbert 对应。

\begin{thm}\label{thm:1.5}
（完整而精确的陈述见定理 3.9。）设 $X$ 是 $k$ 上的光滑刚性解析簇。则
\begin{enumerate}
\item $D^i_{\dR}$ 是从 $X_{\text{\'et}}$ 上的 $\Qp$ 局部系统范畴到 $X$ 上带可积联络的向量丛范畴的函子。
若 $X$ 是一点，则 $D^0_{\dR}$ 与 Fontaine 的 $D_{\dR}$ 函子一致。
\item 函子 $D^i_{\dR}$ 与任意拉回交换。
\item 若 $X$ 连通，且存在经典点 $x$ 使 $\bar L_x$ 为 de Rham，
则 $D^0_{\dR}(L)$ 上存在由子丛给出的递减滤过 $\Fil$，
使得 $\bigl(D^0_{\dR}(L),\nabla_L,\Fil\bigr)$ 是 \cite[Definition 7.4]{Sch2} 意义下
与 $L$ 相伴的带可积联络的滤过 $\OX$-模。
换言之，$L$ 是 \cite[Definition 8.3]{Sch2} 意义下的 de Rham 局部系统。
\end{enumerate}
\end{thm}

\begin{rmk}\label{rmk:1.8}
\begin{enumerate}
\item 尽管对大多数应用（例如定理 \ref{thm:1.3}）只需用到 $D^0_{\dR}$，
但为了证明上述关于 $D^0_{\dR}$ 的陈述，我们需要同时研究所有的 $D^i_{\dR}$。
\item $D^0_{\dR}$ 的局部版本本质上已出现在 Brinon 的工作 \cite{Br2} 中。
\end{enumerate}
\end{rmk}

\begin{rmk}\label{rmk:1.9}
\begin{enumerate}
\item 我们要强调：即使 $(X,L)$ 是 $k$ 上某个代数簇上平展局部系统的解析化，
上述定理中的向量丛 $D^0_{\dR}(L)$ 先验地也是解析的。
这与经典情形一致：在经典情形，Riemann--Hilbert 对应首先被建立为
Betti 局部系统与带平坦联络的复向量丛之间的等价。
不过我们预期 $D^0_{\dR}(L)$ 可以延拓为某个良好紧化 $\bar X$ 上的对数联络，
从而是具有正则奇点的代数联络。
\item 注意向量丛 $D^i_{\dR}(L)$ 一般而言未必具有正确的秩。
因此算术 Riemann--Hilbert 对应只对 de Rham 局部系统表现良好。
事实上，若 $X$ 是（未必光滑的）代数簇，则存在一个定义良好的 de Rham 反常层范畴 $\PerDR(X)$；
我们希望把 $D^0_{\dR}$ 延拓为从 $\PerDR(X)$ 到 $X$ 上代数 $D$-模范畴的函子。
我们计划在后续工作中研究这些延拓。
\end{enumerate}
\end{rmk}

下面解释函子 $\mathrm{RH}$ 的构造思路。它基于相对 $p$ 进 Hodge 理论的近期进展
\cite{Sch1,Sch2,KL1,KL2}。我们将沿用 \cite{Sch2} 的记号。

首先回忆，在经典 Riemann--Hilbert 对应中，从复解析簇 $X$ 上的局部系统
到带联络的向量丛的函子由与解析函数层作张量积给出：
\[
  L\longmapsto L\otimes_{\C}\OX,\qquad \nabla=1\otimes d.
\]
在 $p$ 进情形，我们的函子 $\mathrm{RH}$ 在性质上相同，只不过众所周知，
$p$ 进情形下的解析拓扑（甚至平展拓扑）不够精细，$X$ 上的解析函数层无法胜任。
于是思路是：考虑 $X$ 上某个更精细的拓扑，以及该拓扑下的某个周期层，作为 $\OX$ 的替代。
更精确地说，Scholze \cite{Sch2} 引入了 $X$ 的 pro-平展 site $X_{\proet}$，
作为 $X$ 上通常平展拓扑的加细，因此存在自然投影 $\nu:X_{\proet}\to X_{\text{\'et}}$。
设 $\nu':X_{\proet}/\Xk\to(\Xk)_{\text{\'et}}$ 是 $\nu$ 的限制。
于是每个 $\Qp$ 局部系统 $L$ 都（大致通过拉回）给出 $X_{\proet}$ 上一个局部常值的 $\Qp$-模层，
记作 $\widehat L$。此外，$X_{\proet}$ 上有 de Rham 周期层 $\mathcal O_{B_{\dR}}$，
它是结构层上的一个模，带有可积联络和满足 Griffiths 横截性的递减滤过。然后定义
\[
  \mathrm{RH}(L):=R\nu'_*\bigl(\widehat L\widehat{\otimes}_{\Qp}\mathcal O_{B_{\dR}}\bigr).
\]
由定义可知，$\mathcal O_{B_{\dR}}$ 上的联络诱导了 $\mathrm{RH}(L)$ 上的可积联络 $\nabla_L$。
关键在于证明 $\mathrm{RH}(L)$——它先验地可能是一个层的复形——
实际上是有限生成的 $\OX\widehat{\otimes}\BdR$-模。

为证明这一点，我们先把 $\mathcal O_{B_{\dR}}$ 换成它的零次分次片
$\mathcal O_C=\gr^0\mathcal O_{B_{\dR}}$。于是需要研究
\[
  H(L):=R\nu'_*\bigl(\widehat L\widehat{\otimes}_{\Qp}\mathcal O_C\bigr).
\]
注意：对 $\mathcal O_{B_{\dR}}$ 上联络取分次，定义了 $\mathcal O_C$ 上的一个 Higgs 场，
从而定义了 $H(L)$ 上的一个 Higgs 场 $\vartheta_L$。结果表明，函子
\[
  H:L\longmapsto\bigl(H(L),\vartheta_L\bigr)
\]
正是（$p$ 进 Simpson 对应的一个特例）$p$ 进 Simpson 对应；
后者最早由 Faltings \cite{Fa} 提出，近来由 Abbes--Gros 与 Tsuji \cite{AG,AGT}
系统地研究。\footnote{%
Abbes 指出，将我们的构造与 Abbes--Gros 在 \cite{AGT} 中的构造作比较存在微妙之处，
尽管我们相信二者本质上是相同的。}
我们注意到，这些工作研究的是比本文所考虑的情形更一般、更精细的理论；
但在我们的特殊情形，情况更简单也更漂亮。例如，我们有下述陈述。

\begin{thm}\label{thm:1.6}
（完整而精确的陈述见定理 2.1。）设 $X$ 是 $k$ 上的光滑刚性解析簇。
则 $H$ 是从 $X_{\text{\'et}}$ 上的 $\Qp$ 局部系统范畴到 $\Xk$ 上幂零 Higgs 丛范畴的张量函子。
函子 $H$ 与任意态射的拉回相容，并且在某些条件下与光滑紧合态射的顺像相容。
\end{thm}

\begin{rmk}\label{rmk:1.10}
与 \cite{Fa,AGT} 中的 $p$ 进 Simpson 对应相比，我们的定理似乎包含下述额外结果。
\begin{enumerate}
\item 我们证明 $H(L)$ 上的 Higgs 场是幂零的。这反映了 $L$ 是定义在 $X$ 上
（而非 $\Xk$ 上）这一事实，因而它获得了 $\Gal(\bar k/k)$ 作用带来的额外对称性。
这与 Simpson 的下述结果 \cite[Corollary 4.2]{Sim} 类似：
若 $L$ 是复光滑射影簇上的 Betti 局部系统，且它来自一个 Hodge 结构的变分，
则在（经典）Simpson 对应下相应的 Higgs 丛 $(M,\vartheta)$ 是幂零的。
事实上，$k$ 的分圆塔的 Galois 群应当被视为 Hodge 环面 $\C^\times$ 的 $p$ 进对应物
（见引理 2.15 的证明）。
\item 我们建立了该函子的一些函子性，例如它与拉回的相容性
（这是定理 \ref{thm:1.3} 证明中的一个重要环节）。
\end{enumerate}
\end{rmk}

下面简述全文的组织。本文实际上按与引言相反的顺序写成：
我们先在第 2 节讨论 $p$ 进 Simpson 对应，
然后第 3 节研究 $p$ 进 Riemann--Hilbert 对应，
最后第 4 节证明定理 \ref{thm:1.1} 并给出在志村簇上的应用。

\subsection*{记号与约定}
\addcontentsline{toc}{subsection}{记号与约定}

本文中，$k$ 表示一个 $p$ 进域，$\khb$ 表示 $k$ 的完备化代数闭包。
设 $k_\infty=\bigcup_m k_m$ 是 $k$ 在 $\bar k$ 中的分圆扩张。
设 $K\subset\khb$ 是包含 $k_\infty$ 的完美胚域，$K^\flat$ 表示 $K$ 的倾斜。
我们固定 $k_\infty$ 中一族相容的 $p$ 的幂次单位根 $\{\zeta_{p^m}\}_{m\ge 0}$，
它给出 $\epsilon=(1,\zeta_p,\dots)\in\mathcal O_{K^\flat}$。

我们把 $L$ 上（$L=k$ 或 $K$）的刚性解析簇视为局部有限型的 adic 空间。
特别地，仿射体（affinoid）写作 $\Spa(A,A^+)$，其中 $A^+=A^\circ$。
若 $X=\Spa(A,A^+)$ 是 $k$ 上的仿射体，则记 $A_K=A\widehat{\otimes}_k K$，
$X_K=\Spa(A_K,A_K^+)$ 为其到 $K$ 的基变换。
仿射体之间的平展态射若可写成有理局部化与有限平展态射的复合，则称为标准平展的。
从 $L$ 上刚性解析簇 $X$ 到环面
\[
  T^n=\Spa\bigl(L\langle T_1^{\pm 1},\dots,T_n^{\pm 1}\rangle,\,
  \mathcal O_L\langle T_1^{\pm 1},\dots,T_n^{\pm 1}\rangle\bigr)
\]
的标准平展态射，称为 $X$ 的一个环面坐标卡（toric chart）。

设 $X$ 是 $k$ 上的刚性解析簇。沿用 \cite{Sch2}，记 $X_{\proet}$ 为 $X$ 上的 pro-平展 site，
并记
\[
  \nu:X_{\proet}\to X_{\text{\'et}}
\]
为自然投影。设 $\widehat{\Z}_p$（相应地 $\widehat{\Q}_p$）为 $X_{\proet}$ 上与
$\Z_p$（相应地 $\Qp$）相伴的常值层；对 $X_{\text{\'et}}$ 上的 $\Z_p$ 局部系统 $L$
（相应地 $\Qp$ 局部系统），记 $\widehat L$ 为 $X_{\proet}$ 上与 $L$ 相伴的
$\widehat{\Z}_p$-模（相应地 $\widehat{\Q}_p$-模）
（参见 \cite[§8.2]{Sch2}、\cite[§9.1]{KL1}）。

\subsection*{致谢}
\addcontentsline{toc}{subsection}{致谢}

作者感谢 A. Abbes、P. Colmez、K. S. Kedlaya、蓝凯文（K.-W. Lan）、P. Scholze
与田野的有益讨论，感谢清水浩二（Koji Shimizu）非常仔细地阅读本文并提出宝贵意见。
部分工作完成时，第一作者正在加州理工学院访问，第二作者正在北京国际数学研究中心访问；
作者感谢这些机构的款待。
